% GENERATED FILE. Edit canonical lessons, then run npm run generate:books.
% canonical-source-hash: fnv1a64:e45a7b347790fb7c
% canonical-lessons: KA-C46-fruit, KA-C46-cloth, KA-C46-lamp, KA-C46-salt, KA-C46-book

\chapter{Things You Ask For}
\label{ch:ka-ask-for}
\hlchaptermodality{46}

\begin{chapteropening}
\textbf{By the end of this chapter:} \emph{I can name five everyday things I might ask someone for in Kannada, and say which of them Kannada has always had and which arrived from elsewhere.}

Complete KA-C46-book and say all five words in order.
\end{chapteropening}

\section[haṇṇu]{\kn{ಹಣ್ಣು} (haṇṇu) — a fruit}

\label{lesson:KA-C46-fruit}



\begin{quote}
Before the new run of words: what did \emph{heṇḍati} mean?
\end{quote}

\subsection*{You'll want to know: \kn{ಹಣ್ಣು}}
\textbf{\kn{ಹಣ್ಣು}} (\emph{haṇṇu}) — \textquotedblleft{}a fruit\textquotedblright{}.

Behind it there is a verb: \kn{ಹಣ್ಣಾಗು}, to ripen. The fruit is named for what happened to it — the ripened thing — and Kannada saw no need to invent a second noun for the object itself.

Proto-Dravidian \textbf{*paḻ}, 'to ripen, to grow old', is the source, and the sound law this book keeps returning to is at work in the very first letter. Tamil kept \ta{பழம்} (\emph{paḻam}); Telugu has \te{పండు} (\emph{paṇḍu}); Kannada turned that \emph{p-} into \emph{h-} and arrived at \kn{ಹಣ್ಣು}. It is the same law that made milk \kn{ಹಾಲು} here and left it \emph{pāl} next door.

\begin{grammarlens}[title={What you've built}]
The first of five things you might ask someone for.
\end{grammarlens}

\subsection*{Your turn}
Say these aloud:

\begin{itemize}
\raggedright
  \item \emph{haṇṇu}
  \item it once more, slowly
  \item \emph{haṇṇu}, then \emph{dayaviṭṭu haṇṇu koḍi} — \textquotedblleft{}please give a fruit\textquotedblright{}
\end{itemize}

\begin{tcolorbox}[breakable,colback=teal!4,colframe=teal!35!black,title={Before you move on}]
What does \kn{ಹಣ್ಣು} mean? (\textquotedblleft{}a fruit\textquotedblright{}.) Say it once more.
\end{tcolorbox}
\section[baṭṭe]{\kn{ಬಟ್ಟೆ} (baṭṭe) — cloth}

\label{lesson:KA-C46-cloth}



\begin{quote}
Before the new one: what did \emph{haṇṇu} mean?
\end{quote}

\subsection*{You'll want to know: \kn{ಬಟ್ಟೆ}}
\textbf{\kn{ಬಟ್ಟೆ}} (\emph{baṭṭe}) — \textquotedblleft{}cloth\textquotedblright{}.

\kn{ಬಟ್ಟೆ} is the plain word for cloth: the bolt on the shelf, the washing on the line, the rag in your hand.

Sanskrit has \emph{paṭṭa}, 'a strip of cloth, a bandage', and here the borrowing most likely ran the unusual way — out of Dravidian and into Sanskrit, rather than the direction this book has shown so often. The same strip reached English by another road: Hindi \emph{paṭṭī}, a wrapping, is the soldier's \emph{puttee}.

Older Kannada also used \kn{ಬಟ್ಟೆ} for a road — the way you go — a sense the modern language has quietly let go of.

\begin{grammarlens}[title={What you've built}]
Two: a fruit and a cloth.
\end{grammarlens}

\subsection*{Your turn}
Say these aloud:

\begin{itemize}
\raggedright
  \item \emph{baṭṭe}
  \item it once more, slowly
  \item \emph{haṇṇu}, then \emph{baṭṭe}, and ask for each one with \emph{dayaviṭṭu}
\end{itemize}

\begin{tcolorbox}[breakable,colback=teal!4,colframe=teal!35!black,title={Before you move on}]
What does \kn{ಬಟ್ಟೆ} mean? (\textquotedblleft{}cloth\textquotedblright{}.) Say it once more.
\end{tcolorbox}
\section[dīpa]{\kn{ದೀಪ} (dīpa) — a lamp}

\label{lesson:KA-C46-lamp}



\begin{quote}
Before the new one: what did \emph{baṭṭe} mean?
\end{quote}

\subsection*{You'll want to know: \kn{ದೀಪ}}
\textbf{\kn{ದೀಪ}} (\emph{dīpa}) — \textquotedblleft{}a lamp\textquotedblright{}.

Sanskrit \emph{dīpa}, from a root meaning to blaze, to shine. Kannada took it whole and wore not one letter off it.

Put it beside \emph{āvali}, 'a row', and you have \emph{Dīpāvali} — a row of lamps — a festival whose name explains itself the moment the seam is visible.

Kannada has a word of its own as well, \kn{ಹಣತೆ} (\emph{haṇate}), and it means something narrower: the small clay saucer with a wick lying in it. The borrowed word covers every lamp; the native one names a particular thing you can hold in one hand.

\begin{grammarlens}[title={What you've built}]
Three, and one of them came from Sanskrit.
\end{grammarlens}

\subsection*{Your turn}
Say these aloud:

\begin{itemize}
\raggedright
  \item \emph{dīpa}
  \item it once more, slowly
  \item all three so far, and say which of them a potter makes
\end{itemize}

\begin{tcolorbox}[breakable,colback=teal!4,colframe=teal!35!black,title={Before you move on}]
What does \kn{ದೀಪ} mean? (\textquotedblleft{}a lamp\textquotedblright{}.) Say it once more.
\end{tcolorbox}
\section[uppu]{\kn{ಉಪ್ಪು} (uppu) — salt}

\label{lesson:KA-C46-salt}



\begin{quote}
Before the new one: what did \emph{dīpa} mean?
\end{quote}

\subsection*{You'll want to know: \kn{ಉಪ್ಪು}}
\textbf{\kn{ಉಪ್ಪು}} (\emph{uppu}) — \textquotedblleft{}salt\textquotedblright{}.

Proto-Dravidian \textbf{*uppu}, and the whole family kept it very nearly untouched: Tamil \ta{உப்பு}, Telugu \te{ఉప్పు}, Malayalam \ml{ഉപ്പ്}.

There is a \emph{p} sitting in the middle of this word, and it did not become \emph{h}. That is worth stopping over, because the law that made \kn{ಹಣ್ಣು} out of an older \emph{p}-word only ever reached the FIRST sound of a word. Inside one, the \emph{p} stays where it was put.

\begin{grammarlens}[title={What you've built}]
Four. Salt, and the fruit, and the cloth, and the lamp.
\end{grammarlens}

\subsection*{Your turn}
Say these aloud:

\begin{itemize}
\raggedright
  \item \emph{uppu}
  \item it once more, slowly
  \item \emph{uppu}, then \emph{haṇṇu}, and hear the \emph{p} survive in one and vanish from the other
\end{itemize}

\begin{tcolorbox}[breakable,colback=teal!4,colframe=teal!35!black,title={Before you move on}]
What does \kn{ಉಪ್ಪು} mean? (\textquotedblleft{}salt\textquotedblright{}.) Say it once more.
\end{tcolorbox}
\section[pustaka]{\kn{ಪುಸ್ತಕ} (pustaka) — a book}

\label{lesson:KA-C46-book}



\begin{quote}
Before the new one: what did \emph{uppu} mean?
\end{quote}

\subsection*{You'll want to know: \kn{ಪುಸ್ತಕ}}
\textbf{\kn{ಪುಸ್ತಕ}} (\emph{pustaka}) — \textquotedblleft{}a book\textquotedblright{}.

Sanskrit \emph{pustaka} is the source, but the trail behind it leaves India altogether. Middle Persian \emph{post} meant 'hide, skin' — the thing written on before there was paper.

Now look at its first letter. \kn{ಪುಸ್ತಕ} opens with \emph{p-}, and that \emph{p-} stayed. So did the \emph{p-} of every other word Kannada took in from Sanskrit. The law that turned inherited \emph{p-} into \emph{h-} had finished its work before these words walked in, and it left them alone.

Two of these five came from outside; three were here already.

\begin{grammarlens}[title={What you've built}]
Five things you might ask for: a fruit, a cloth, a lamp, salt, a book.
\end{grammarlens}

\subsection*{Your turn}
Say these aloud:

\begin{itemize}
\raggedright
  \item \emph{pustaka}
  \item it once more, slowly
  \item all five in order, then sort them into the inherited and the borrowed
\end{itemize}

\begin{tcolorbox}[breakable,colback=teal!4,colframe=teal!35!black,title={Before you move on}]
What does \kn{ಪುಸ್ತಕ} mean? (\textquotedblleft{}a book\textquotedblright{}.) Say it once more.
\end{tcolorbox}
